\documentclass{ximera}
%geschreven 16/11/26
\title{Spanningsbronnen}
\author{Daan Lipkens}

\begin{document}
\begin{abstract}
    Gelijkspanningsbronnen en wisselspanningsbronnen en hun symbolen in een elektrisch schema.
\end{abstract}
\maketitle
\label{app:spanningsbronnen}

\section{Spanningsbronnen}\label{sec:spanningsbronnen}
Een \textbf{spanningsbron} (bijvoorbeeld een batterij) onderhoudt de spanning tussen haar twee polen. Bij een \textbf{gelijkspanningsbron} is deze spanning constant in de tijd. De pluspool is de pool met de hoogste potentiaal. Bij een \textbf{wisselspanningsbron} (zoals de netspanning) wisselen de polen voortdurend van teken. In Europa doorloopt de spanning $50$ volledige wisselingen per seconde ($50 \, \mathrm{Hz}$). Een gelijkspanningsbron houdt het ladingsverschil tussen de polen in stand: de minpool heeft een overschot aan elektronen, de pluspool een tekort. Hoe groter dat verschil, hoe groter de spanning.

\begin{remark}[expandable=true,title={Herhaling: symbolen van spanningsbronnen}]\nl
In een elektrisch schema krijgt elke spanningsbron een eigen symbool. Bij een gelijkspanningsbron staat de lange streep aan de pluspool en de korte streep aan de minpool. Een variabele gelijkspanningsbron herken je aan de schuine pijl door het symbool: de spanning kan je regelen. Het symbool van een wisselspanningsbron is een cirkel met een sinusvormig golfje, omdat de spanning voortdurend van teken wisselt.
\begin{center}
    \begin{tikzpicture}[scale=1]
        % Gelijkspanningsbron: lange streep = pluspool, korte streep = minpool
        \begin{scope}
            \draw[thick] (-0.7,0) -- (-0.1,0);
            \draw[line width=0.4pt] (-0.1,-0.45) -- (-0.1,0.45);
            \draw[line width=2pt] (0.1,-0.22) -- (0.1,0.22);
            \draw[thick] (0.1,0) -- (0.7,0);
            \node[above left] at (-0.1,0.45) {$+$};
            \node[above right] at (0.1,0.22) {$-$};
            \node[below] at (0,-0.7) {gelijkspanningsbron};
        \end{scope}
        % Variabele gelijkspanningsbron: zelfde symbool met een schuine pijl van linksonder naar rechtsboven
        \begin{scope}[xshift=4.2cm]
            \draw[thick] (-0.7,0) -- (-0.1,0);
            \draw[line width=0.4pt] (-0.1,-0.45) -- (-0.1,0.45);
            \draw[line width=2pt] (0.1,-0.22) -- (0.1,0.22);
            \draw[thick] (0.1,0) -- (0.7,0);
            \draw[-{Latex[length=2.5mm]}, thick] (-0.5,-0.55) -- (0.5,0.55);
            \node[above left] at (-0.1,0.45) {$+$};
            \node[right] at (0.15,-0.3) {$-$};
            \node[below] at (0,-0.7) {variabele gelijkspanningsbron};
        \end{scope}
        % Wisselspanningsbron: cirkel met een sinusvormig golfje
        \begin{scope}[xshift=8.4cm]
            \draw[thick] (0,0) circle (0.45);
            \draw[thick, domain=-0.3:0.3, samples=40, smooth] plot (\x,{0.22*sin(\x*360/0.6)});
            \node[below] at (0,-0.7) {wisselspanningsbron};
        \end{scope}
    \end{tikzpicture}
\end{center}
\end{remark}

\end{document}
